\documentclass[12pt]{article}
\usepackage[utf8]{inputenc}
\usepackage[T1]{fontenc}
\usepackage[margin=1in]{geometry}
\usepackage{setspace}
\usepackage{xcolor}
\usepackage[colorlinks=true, urlcolor=blue!60!black, linkcolor=blue!60!black]{hyperref}
\setlength{\parindent}{2em}
\setlength{\parskip}{1.25em}
\renewcommand{\baselinestretch}{1.0}
\begin{document}
\begin{center}
{\LARGE \textit{Boumediene v. Bush}, 553 U.S. 723 (2008)\par}
\vspace{0.8em}
{\large Case Brief --- Four-Page Edition\par}
\end{center}
\vspace{1em}

\section{Memory Jogger}

Holding 5--4 that the Suspension Clause has full effect at Guantanamo Bay, the Court held \S~7 of the Military Commissions Act unconstitutional as applied to these detainees, because the Detainee Treatment Act's review of status determinations was not an adequate and effective substitute for habeas corpus (\href{https://supreme.justia.com/cases/federal/us/553/723/}{553 U.S. 723}).

\section{Detailed Case Facts}

After the September 11 attacks, Congress passed the Authorization for Use of Military Force. Lakhdar Boumediene, an Algerian-born Bosnian citizen seized in Bosnia, and other foreign nationals captured abroad were held at the U.S. Naval Station at Guantanamo Bay, where Cuba kept technical \textit{de jure} sovereignty under a 1903 lease while the United States exercised complete and indefinite \textit{de facto} control. Each denied being an enemy combatant.

After \textit{Rasul v. Bush} (2004) extended statutory habeas to Guantanamo, the Defense Department created Combatant Status Review Tribunals (CSRTs), at which detainees had no lawyer, only a non-lawyer ``Personal Representative,'' no access to classified evidence, a presumption favoring the Government, and no hearsay bar. The Detainee Treatment Act of 2005 (DTA) channeled limited review of CSRT rulings to the D.C. Circuit. After \textit{Hamdan v. Rumsfeld} (2006), the Military Commissions Act of 2006 (MCA) \S~7 stripped habeas jurisdiction over ``all cases, without exception,'' brought by noncitizens detained as enemy combatants.

\section{Procedural History}

The D.C. Circuit held that MCA \S~7 stripped jurisdiction over the pending cases and that noncitizens held outside sovereign U.S. territory had no constitutional habeas right. After first denying certiorari, the Supreme Court granted review on June 29, 2007, heard argument on December 5, 2007, and decided the case on June 12, 2008.

\section{Judicial Votes}

Reversed and remanded, \textbf{5--4}.
\begin{itemize}
\item \textbf{Majority:} Justice Kennedy, joined by Justices Stevens, Souter, Ginsburg, and Breyer.
\item \textbf{Concurrence:} Justice Souter, joined by Justices Ginsburg and Breyer.
\item \textbf{Dissents:} Chief Justice Roberts, joined by Justices Scalia, Thomas, and Alito; Justice Scalia, joined by the Chief Justice and Justices Thomas and Alito.
\end{itemize}

\section{Holding}

MCA \S~7 strips statutory habeas jurisdiction over pending Guantanamo petitions. The Suspension Clause has full effect at Guantanamo under a functional test, and noncitizen status and enemy-combatant designation do not bar petitioners from invoking it. Because DTA review of CSRT findings is not an adequate and effective substitute, and Congress did not formally suspend the writ, \S~7 is an unconstitutional suspension.

\section{Analysis of Opinions}

\textbf{Majority (Kennedy).} Tracing the writ from Magna Carta, Kennedy treated the Suspension Clause as a structural safeguard against executive detention. Finding the 1789 record inconclusive, he rejected \textit{de jure} sovereignty as the boundary, since leasing foreign territory would otherwise let the political branches evade the Constitution: ``questions of extraterritoriality turn on objective factors and practical concerns, not formalism.'' Three factors govern: the detainees' status and the adequacy of the process that determined it (far weaker than the trials in \textit{Johnson v. Eisentrager}); the sites of capture and detention (a century of exclusive U.S. control); and practical obstacles (none, since Guantanamo is no war zone). An adequate substitute must let a court correct errors, assess the evidence, consider newly discovered exculpatory evidence, and order conditional release. The DTA did not, and after six years of detention, exhaustion was not required.

\textbf{Souter (concurring).} \textit{Rasul} had already recognized U.S. control of Guantanamo, and six years without meaningful review answered the charge of prematurity.

\textbf{Roberts (dissenting).} The Court struck down ``the most generous set of procedural protections ever afforded aliens detained by this country as enemy combatants'' without letting the D.C. Circuit test the DTA.

\textbf{Scalia (dissenting).} \textit{Eisentrager} barred habeas for enemy aliens held abroad, the 1789 writ never ran to non-subjects abroad, and the ``judicially brainstormed'' test usurps war powers.

\textbf{Later cases.} \textit{Al Maqaleh v. Gates} (D.C. Cir. 2010) held the Clause inapplicable at Bagram, and \textit{DHS v. Thuraissigiam} (2020) limited it to the writ ``as it existed in 1789.''

\section{Examples: Future Applications}

\begin{enumerate}
\item \textbf{Same-Side (Applies): Leased Atoll.} A remote atoll under a 99-year lease giving exclusive U.S. control, far from any war, with a review board barring counsel, triggers the Clause.
\item \textbf{Same-Side (Inadequate): No Release Power.} A new statute gives detainees counsel, access to the evidence, and the right to present new exculpatory evidence, but lets the court only recommend release to the Secretary of Defense. Without the power to order release, the substitute fails.
\item \textbf{Opposite-Side (Does Not Apply): Forward Base.} A detainee at a temporary base in an active combat zone likely has no constitutional habeas right.
\item \textbf{Opposite-Side (Outside the Writ): Expedited Removal.} A noncitizen seeking a new asylum hearing rather than release cannot invoke \textit{Boumediene}.
\item \textbf{Fence-Sitter: Allied Base.} A U.S.-run facility on an ally's base, far from hostilities, has weak process and no war friction but lacks Guantanamo's complete control, and a court order might create friction with the host government.
\end{enumerate}

\section{Critique}

\textbf{Indeterminacy and missing rules.} The functional test prevents formalistic evasion but gives little guidance on when control becomes \textit{de facto} sovereignty. The Court also left open the definition of ``enemy combatant,'' the burden of proof, and evidentiary rules, inviting years of D.C. Circuit litigation that narrowed detainee victories.

\textbf{Tension with \textit{Thuraissigiam}.} \textit{Boumediene}'s functionalism, which found the 1789 record inconclusive, contrasts with \textit{Thuraissigiam}'s strict 1789 originalism, a shift reflecting the Court's changed composition.

\section{Key Quotations}

\begin{quote}
``To hold the political branches have the power to switch the Constitution on or off at will is quite another.''
\end{quote}
\noindent --- Justice Kennedy, for the Court, 553 U.S. at 765.

\begin{quote}
``However, security subsists, too, in fidelity to freedom's first principles, chief among them being freedom from arbitrary and unlawful restraint and the personal liberty that is secured by adherence to the separation of powers.''
\end{quote}
\noindent --- Justice Kennedy, for the Court, 553 U.S. at 797.

\begin{quote}
``Today, for the first time in our Nation's history, the Court confers a constitutional right to habeas corpus on alien enemies detained abroad by our military forces in the course of an ongoing war. \ldots{} The Nation will live to regret what the Court has done today.''
\end{quote}
\noindent --- Justice Scalia, dissenting, 553 U.S. at 826--27, 850.

\end{document}
